%%%%%%%%%%%%%%%%%%%%%%%%%%%%%%%%%%%%%%%%%%%%%%%%%%%%%%%%%%%%%%%%%%%%%%%%%%%%%%%%
% ABSTRACT

\begin{abstract}

%% neste arquivo abstract.tex
%% o texto do resumo e as palavras-chave têm que ser em Inglês para os documentos escritos em Português
%% o texto do resumo e as palavras-chave têm que ser em Português para os documentos escritos em Inglês
%% os nomes dos comandos \begin{abstract}, \end{abstract}, \keywords e \palavrachave não devem ser alterados

%\selectlanguage{portuguese}	%% para os documentos escritos em Português
\selectlanguage{portuguese}	%% para os documentos escritos em Inglês

%\hypertarget{estilo:abstract}{} %% uso para este Guia

\textcolor{red}{Uma nova parametrização de cold pool, acoplada com a parametrização de convecção Grell-Freitas, foi avaliada no mais novo sistema brasileiro de previsão do tempo e clima, o \textit{Model for Ocean-laNd-Atmosphere predictioN} (MONAN), cujo núcleo dinâmico é baseado no \textit{Model for Prediction Across Scales} (MPAS). Essa iniciativa liderada pelo INPE visa desenvolver um sistema de modelagem unificado capaz de produzir previsões ao longo de diversas escalas temporais e espaciais, melhorando a previsão do tempo e clima para o Brasil e a América do Sul. O estudo realizado nessa dissertação avalia o impacto da parametrização de \textit{cold pool} na trajetória, intensidade e chuva associadas à dois ciclones tropicais da temporada de 2024 na bacia do Atlântico Norte: Beryl e Helene. Ciclones tropicais são sistemas de baixa pressão intensos e organizados, associados a ventos fortes, ressacas e inundações, o que torna a sua previsão extremamente importante para a sociedade. Para o Furacão Beryl, experimentos de sensibilidade foram realizados afim de avaliar as influências nas condições iniciais, no tempo de vida das \textit{cold pools}, na distância vertical em que o fluxo de massa ocorre e na resolução horizontal (15, 30 e 60 km). A inclusão da dinâmica das \textit{cold pools} melhoraram as previsões das trajetórias, particularmente para maiores tempos de previsão; e na representação da chuva, especialmente na organização espacial das espirais das bandas de chuva e das chuvas intensas. A previsão da intensidade do vento também foi melhorada por conta do benefício de utilizar uma representação mais realista da estrutura convectiva, apesar de um viés sistemático na pressão. Uma configuração ótima, combinando a redução do tempo de vida das \textit{cold pools} (1 hora), a altura deslocada do pico do fluxo de massa e o aumento da resolução, resultou em um desempenho equilibrado em todas as métricas avaliadas. Essa configuação foi sobretudo efetiva na representação das estruturas da chuva, particularmente nos acumulados de precipitação, uma característica dos impactos dos ciclones tropicais. Quando aplicada ao Furacão Helene, essa configuração foi capaz de reproduzir com maior fidelidade o evento extremo de precipitação no sudeste dos Estados Unidos, apresentando melhorias na distribuição espacial, na intensidade e na acurácia da trajetória. Esses resultados ressaltam o papel substancial dos parâmetros das \textit{cold pools} nas simulações de ciclones tropicais no MONAN, particularmente na previsão de chuva.}

\keywords{%
	\palavrachave{Cold Pool}%
	\palavrachave{Parametrização}%
	\palavrachave{MONAN}%
	\palavrachave{Furacões}%
	\palavrachave{Beryl}%
	\palavrachave{Helene}
}

%\selectlanguage{portuguese}	%% para os documentos escritos em Português
\selectlanguage{english}	%% para os documentos escritos em Inglês

\end{abstract}